\documentclass[10pt,a4paper]{amsart}
\usepackage[margin=19mm]{geometry}
\usepackage{amsmath,amssymb,mathtools}
\usepackage[scr=rsfs,cal=euler]{mathalfa}
\usepackage{xfrac}
\usepackage[unicode,hidelinks,bookmarks=false]{hyperref}
\newcommand{\ie}{i.e.\ }
\newcommand{\cf}{cf.\ }
\newcommand{\resp}{respectively\ }
\newcommand{\Subsection}{\subsection}
\providecommand{\nobreakdash}{\nobreak\hbox{-}}
\setlength{\emergencystretch}{2em}
\multlinegap0pt
\usepackage{amssymb}
\newtheorem{theorem}{Theorem}
\title{third-order relativistic kinematics and siacci shock decompositions via localized darboux frames}
\author{Fatma ALMAZ, \c{C}\.{ı}\u{g}dem YAVUZ}
\date{Source: arXiv:2609.05497v1 (CC BY 4.0)}
\begin{document}
\maketitle

\noindent\emph{Source and licence.} third-order relativistic kinematics and siacci shock decompositions via localized darboux frames. Version arXiv:2609.05497v1 (CC BY 4.0), distributed by arXiv under the Creative Commons Attribution 4.0 International licence, http://creativecommons.org/licenses/by/4.0/, as recorded in the arXiv licence metadata for this exact version and shown on the abstract page https://arxiv.org/abs/2609.05497v1. Full licence notice: \texttt{licenses/arxiv-2609.05497-CC-BY-4.0.txt}.\par\smallskip

\noindent\textit{Editorial scope.} Counted unit: the article's theorem theorem (source0.tex lines 477--500). The article's complete original proof of it is delivered with it (source0.tex lines 502--591, one \texttt{proof} environment of 90 lines). No mathematics is written, repaired or re-derived: the statement and the proof are the article's own lines, in the article's own notation, and no retained line is substituted or re-flowed.  No further source-local context is needed: every cross-reference of every retained line resolves inside the two delivered spans.  Nothing was omitted from any retained span, so the package carries no omission record.  No retained line contains a citation, so the wrapper carries no bibliography and no bibliography entry.  No retained line contains a figure environment, an image-inclusion command or drawing code, so no image is delivered and the package declares no asset dependency.  Editorial formatting only, in addition: an AMS \texttt{amsart} preamble that restates the article's own declarations and macros used by the retained lines, namely the environments \texttt{theorem} on separate counters, together with the standard packages those lines need.  The retained environments print in this document as Theorem 1 in order of appearance, whereas the article prints the same items with the numbering of its own full document.  The complete native source of the article as distributed in the arXiv e-print for this exact version is bundled unchanged as \texttt{sources/209440/source0.tex}.
\begin{theorem}
In Minkowski 3-space $\mathbb{E}_1^3$, consider a physical particle $p$ moving along a regular trajectory $\gamma$ that conforms to a localized Darboux frame. Assume that this particle is strictly constrained to a fixed plane $\text{Span}\{T, \sinh \varphi N_{M} + \cosh \varphi G\}$ or $\text{Span}\{T, \cosh \varphi N_{M} + \sinh \varphi G\}$ that isolates the spatial coordinate origin $O$. Furthermore, assume that the component of the particle's angular momentum vector along the normal vector field of the respective spanning plane is nowhere vanishing. Under these boundary conditions, the relativistic jerk vector $J$ is uniquely expressed as:
\begin{eqnarray*}
J &=&\left( \frac{d^{3}s}{dt^{3}}+\kappa ^{2}(\frac{ds}{dt})^{3}-\frac{%
\omega _{1}}{\omega _{2}}(3\kappa \frac{ds}{dt}\frac{d^{2}s}{dt^{2}}+\kappa
^{\prime }(\frac{ds}{dt})^{3})\right) \overrightarrow{T} \\
&&+\left( \frac{\sqrt{\vert\omega _{1}^{2}-\omega _{2}^{2}\vert}}{\omega _{2}}\left(
3\kappa \frac{ds}{dt}\frac{d^{2}s}{dt^{2}}+\kappa ^{\prime }(\frac{ds}{dt}%
)^{3}\right) \right) \overrightarrow{J_{\lambda }}
\end{eqnarray*}%
yielding a vanishing second radial component ($J^{\zeta} = 0$). Alternatively, if the geometric sub-case satisfies $\kappa = 0 \wedge \varphi^{\prime} + \tau_g = 0$, the multi-directional jerk field collapses entirely to a pure one-dimensional translational thrust wave: 
\begin{equation*}
J = \frac{d^{3}s}{dt^{3}} T, \quad J^{\lambda} = J^{\zeta} = 0.
\end{equation*}

The first radial component $J^{\lambda}$ remains invariant in the first configuration if and only if $\cosh \varphi N_{M}+\sinh \varphi G = \text{constant}$; otherwise, it vanishes identically. Concurrently, the global position vectors of the particle $p$ conform to the following metric restrictions:
\begin{equation*}
\gamma = \sqrt{\left|\omega _{1}^{2}-\omega _{2}^{2}\right|} J_{\lambda} + \omega _{3} \left( \cosh \varphi N_{M} + \sinh \varphi G \right), 
\end{equation*}
and
\begin{equation*}
\gamma = \omega _{2} \left( \sinh \varphi N_{M} + \cosh \varphi G \right) + \sqrt{\omega _{1}^{2} + \omega _{3}^{2}} J_{\zeta}. 
\end{equation*}
\end{theorem}

\begin{proof}
In Minkowski 3-space, let the trajectory of a physical particle $p$ be bounded by a regular curve $\gamma$ lying in a fixed plane $\text{Span}\{T, \sinh \varphi N_{M} + \cosh \varphi G\}$ that isolates the spatial coordinate origin. Under this geometric constraint, the out-of-plane positional projection coefficient $\omega_{3}$ is strictly non-zero. Furthermore, the vector field in the parallel plane $\text{Span}\{T, \cosh \varphi N_{M} + \sinh \varphi G\}$ explicitly defines the normal vector field to the initial constraint plane $\text{Span}\{T, \sinh \varphi N_{M} + \cosh \varphi G\}$, meaning it functions as an invariant directional field along the trajectory $\gamma$. By invoking the constancy of this normal vector field along the worldline, we evaluate its intrinsic derivative with respect to the arc length parameter $s$, yielding the following system:
\begin{equation}
\frac{d}{ds}\left( \cosh \varphi \overrightarrow{N_{M}}+\sinh \varphi 
\overrightarrow{G}\right) =\left( 
\begin{array}{c}
-k_{g}\cosh \varphi \\ 
+k_{n}\sinh \varphi%
\end{array}%
\right) \overrightarrow{T}+(\varphi ^{\prime }+\tau _{g})\left( 
\begin{array}{c}
\sinh \varphi \overrightarrow{N_{M}} \\ 
+\cosh \varphi \overrightarrow{G}%
\end{array}%
\right)  \tag{3.18}
\end{equation}%
similarly, for the alternative configuration where $\omega_{2}$ is non-zero, the vector field spanning $\text{Span}\{T, \sinh \varphi N_{M} + \cosh \varphi G\}$ defines the geometric normal to the spanning plane $\text{Span}\{T, \cosh \varphi N_{M} + \sinh \varphi G\}$. From the spatial invariance of this normal vector field along the trajectory, the structural derivative relation satisfies:
\begin{equation}
\frac{d}{ds}\left( \sinh \varphi \overrightarrow{N_{M}}+\cosh \varphi 
\overrightarrow{G}\right) =\left( 
\begin{array}{c}
k_{n}\cosh \varphi \\ 
-k_{g}\sinh \varphi%
\end{array}%
\right) \overrightarrow{T}+(\varphi ^{\prime }+\tau _{g})\left( 
\begin{array}{c}
\cosh \varphi \overrightarrow{N_{M}} \\ 
+\sinh \varphi \overrightarrow{G}%
\end{array}%
\right) .  \tag{3.19}
\end{equation}

By substituting the localized Darboux identities $k_{g} = \kappa \sinh \varphi$, $k_{n} = \kappa \cosh \varphi$, and $\tau = \tau_{g} + \varphi^{\prime}$ into the derivative expansions (3.18) and (3.19), the systems reduce directly to the following geometric conditions:
\begin{equation}
\frac{d}{ds}\left( \cosh \varphi \overrightarrow{N_{M}}+\sinh \varphi 
\overrightarrow{G}\right) =0.\overrightarrow{T}+(\varphi ^{\prime }+\tau
_{g})\left( \sinh \varphi \overrightarrow{N_{M}}+\cosh \varphi 
\overrightarrow{G}\right) =0  \tag{3.20}
\end{equation}%
and 
\begin{equation}
\frac{d}{ds}\left( \sinh \varphi \overrightarrow{N_{M}}+\cosh \varphi 
\overrightarrow{G}\right) =\kappa \overrightarrow{T}+(\varphi ^{\prime
}+\tau _{g})\left( \cosh \varphi \overrightarrow{N_{M}}+\sinh \varphi 
\overrightarrow{G}\right) =0,  \tag{3.21}
\end{equation}%
and from (3.20) we have 
\begin{equation}
\tau _{g}+\varphi ^{\prime }=0\Rightarrow \varphi =-\int \tau _{g}ds\text{
or }\tau =0,  \tag{3.22}
\end{equation}%
and from (3.21) we get 
\begin{equation}
\kappa \overrightarrow{T}+(\varphi ^{\prime }+\tau _{g})\left( \cosh \varphi 
\overrightarrow{N_{M}}+\sinh \varphi \overrightarrow{G}\right) =0\Rightarrow
\kappa =0\wedge \varphi ^{\prime }+\tau _{g}=0.  \tag{3.23}
\end{equation}

Finally, if the values found in (3.22) are substituted into equations (3.13)
or (3.14), the jerk vector components are obtained as follows%
\begin{equation}
J^{T}=\frac{d^{3}s}{dt^{3}}+\kappa ^{2}\left( \frac{ds}{dt}\right) ^{3}-%
\frac{\omega _{1}}{\omega _{2}}\left( 3\kappa \frac{ds}{dt}\frac{d^{2}s}{%
dt^{2}}+\kappa ^{\prime }\left( \frac{ds}{dt}\right) ^{3}\right)  \tag{3.24}
\end{equation}%
\begin{equation}
J^{\lambda }=\frac{\sqrt{\vert\omega _{1}^{2}-\omega _{2}^{2}\vert}}{\omega _{2}}%
(3\kappa \frac{ds}{dt}\frac{d^{2}s}{dt^{2}}+\kappa ^{\prime }\left( \frac{ds%
}{dt}\right) ^{3});J^{\zeta }=0  \tag{3.25}
\end{equation}%
and if the values {}{}found in (3.23) are substituted into equations (3.13)
or (3.14), the jerk vector components are obtained as follows%
\begin{equation}
J^{T}=\frac{d^{3}s}{dt^{3}};J^{\lambda }=J^{\zeta }=0.  \tag{3.26}
\end{equation}

Furthermore, under these conditions,from (3.12a) and (3.12b) the position
vectors of particle $p$ are written as follows%
\begin{equation}
\gamma =\sqrt{\vert\omega _{1}^{2}-\omega _{2}^{2}\vert}\overrightarrow{J_{\lambda }}%
+\omega _{3}\left( \cosh \varphi \overrightarrow{N_{M}}+\sinh \varphi 
\overrightarrow{G}\right)  \tag{2.27}
\end{equation}%
and%
\begin{equation}
\gamma =\omega _{2}\left( \sinh \varphi \overrightarrow{N_{M}}+\cosh \varphi 
\overrightarrow{G}\right) +\sqrt{\omega _{1}^{2}+\omega _{3}^{2}}%
\overrightarrow{J_{\zeta }}.  \tag{3.28}
\end{equation}
\end{proof}

\end{document}
